% Generated draft: requires mathematical and source review.
\documentclass[11pt]{article}
\usepackage[margin=25mm]{geometry}
\usepackage{fontspec}
\setmainfont{Latin Modern Roman}
\usepackage{amsmath,amssymb,mathtools,unicode-math}
\setmathfont{Latin Modern Math}
\usepackage{enumitem,needspace,xurl,hyperref}
\setlength{\parindent}{0pt}
\setlength{\parskip}{5pt}
\setlength{\emergencystretch}{3em}
\newcommand{\N}{\mathbb{N}}
\newcommand{\Z}{\mathbb{Z}}
\newcommand{\Q}{\mathbb{Q}}
\newcommand{\R}{\mathbb{R}}
\newcommand{\C}{\mathbb{C}}
\newcommand{\F}{\mathbb{F}}
\newcommand{\A}{\mathbb{A}}
\newcommand{\PP}{\mathbb{P}}
\newcommand{\E}{\mathbb{E}}
\newcommand{\eps}{\varepsilon}
\newcommand{\dd}{\,\mathrm{d}}
\begin{document}
\section*{UnsolvedMath ID: 800319}
\textbf{Source problem code:} catalogue entry
\subsection{Definitions and mathematical statement}
Fix an integer \(k\ge 1\). The supplied record does not specify whether graphs are simple, undirected, or allowed to have loops or multiple edges. Under the provisional convention of finite simple undirected graphs, let
\[
 a_{k,n}=\left|\left\{G:\ |V(G)|=n,\ \deg_G(v)=k\text{ for every }v\in V(G)\right\}/\cong\right|,
 \qquad n\ge 1,
\]
where \(\cong\) denotes graph isomorphism. Thus each isomorphism class is counted once, independently of vertex labels. Under this convention, \(a_{k,n}=0\) whenever \(k\ge n\) or \(kn\) is odd.

\textbf{Conjecture 3 (Pak, 2018), provisional precise formulation.} For every fixed integer \(k\ge 1\), there exist a deterministic algorithm \(A_k\), a constant \(C_k>0\), and an integer \(d_k\ge 0\) such that, for every integer \(n\ge 1\), the algorithm \(A_k\), given \(n\), outputs the exact integer \(a_{k,n}\) in binary using at most
\[
 C_k(n+1)^{d_k}
\]
bit operations. The algorithm and polynomial bound may depend on \(k\); no bound uniform in variable \(k\) is asserted. Computing all terms through index \(n\) is an equivalent polynomial-in-\(n\) formulation for fixed \(k\).

The graph convention and the deterministic bit-operation interpretation above are explicit provisional interpretations, not hypotheses or computational conventions established by the supplied excerpt. The record does not specify the index-zero term. Source review is required before treating this as a fully resolved formulation.

The supplied background reports that reference [Ges] establishes similar polynomial recurrences for the numbers of labeled \(k\)-regular graphs for every fixed \(k\ge 1\). This background assertion is distinct from the conjecture about unlabeled graphs; the conjecture does not explicitly demand such a recurrence.
\subsection{Short English statement}
For each fixed degree \(k\ge 1\), Pak conjectures that the number of unlabeled \(k\)-regular graphs on \(n\) vertices can be computed exactly in time polynomial in \(n\). The displayed formulation provisionally uses simple undirected graphs and bit complexity; these conventions require source review. The registry labels the problem open, but that status has not been independently verified.
\subsection{Sources}
\begin{itemize}
\item Local UnsolvedMath registry: permanent ID \(800319\), source problem code \texttt{catalogue entry}, title ``Complexity in Enumerative Combinatorics --- Conjecture 3.'' The supplied record attributes the conjecture to Pak, with proposed year 2018, and labels it \emph{open}; this is supplied registry information, not an independently verified status assessment.
\item Pak, \emph{Complexity Problems in Enumerative Combinatorics} (2018), Conjecture 3. Supplied source URL: \url{https://arxiv.org/abs/1803.06636}. The background records access on 2026-10-16 and extraction by \texttt{arxiv\_latex}; these are supplied provenance details, not a claim of access or verification in this draft.
\item Reference [Ges], cited in the supplied excerpt for polynomial recurrences counting labeled regular graphs. No author, title, publication details, or separate URL for [Ges] is supplied, so this reference cannot be resolved from the record alone.
\end{itemize}
\end{document}
